% =============================================================================
% Method section: Popularity-Aware LightGCN
% -----------------------------------------------------------------------------
% Snippet meant to be \input{} or pasted directly into the report's main .tex
% file at the point where the Method / Approach chapter begins (i.e., nested
% under an existing \section{...} in your document).
%
% Structure: ONE \subsection{Proposed Method} (overview + roadmap) containing
% SIX \subsubsection{...} children, each representing one implemented and
% experimentally validated component.
%
% Required packages (add to preamble if not already present):
%   \usepackage{amsmath, amssymb}
%   \usepackage{algorithm}
%   \usepackage{algorithmic}
%   \usepackage{booktabs}
%   \usepackage{graphicx}   % needed for \resizebox, used by \fitwidthtable below
%
% \fitwidthtable{...} wraps a table and shrinks it ONLY if it is wider than the
% text column (never stretches a narrow table). \providecommand is used so this
% is safe to load even if another \input'd snippet defines it too.
\providecommand{\fitwidthtable}[1]{%
  \setbox0=\hbox{#1}%
  \ifdim\wd0>\linewidth
    \resizebox{\linewidth}{!}{\box0}%
  \else
    \box0
  \fi
}
%
% All formulas below have been verified line-by-line against the actual
% implementation (src/models.py, src/losses.py, src/ile_losses.py,
% src/neg_sampling.py, src/popaware_training.py) as of the final 3-seed
% experiments reported in this work.
%
% NOTE on citation keys: he2020lightgcn, rendle2009bpr, wu2021sgl, and
% wang2020alignuniform refer to real, well-known papers (LightGCN — He et al.,
% SIGIR 2020; BPR — Rendle et al., UAI 2009; SGL — Wu et al., SIGIR 2021;
% alignment/uniformity — Wang & Isola, ICML 2020) but the exact \cite keys must
% be checked against your own .bib file before compiling.
%
% NOTE: "Section~\ref{sec:results}" is a placeholder — replace with \ref{} to the
% actual label of your Experiments section once written.
% =============================================================================

\subsection{Proposed Method}
\label{sec:proposed_method}

Graph-based collaborative filtering models such as LightGCN~\cite{he2020lightgcn} are
particularly susceptible to popularity bias: because the neighborhood-aggregation step
propagates a signal proportional to node degree, high-degree (popular) item nodes receive
disproportionately more gradient updates than low-degree (long-tail) item nodes. Over
training, this drives popular-item embeddings to larger norms and causes them to dominate
the dot-product ranking for almost every user, irrespective of individual user preference.

We propose \emph{Popularity-Aware LightGCN}, which augments the LightGCN backbone with four
complementary interventions. We introduce them in the order below because each one closes a
specific gap left open by the ones before it, rather than as four independent, parallel
tricks; all four are independently switchable, so that disabling all of them recovers plain
LightGCN trained with uniform negative sampling:
\begin{enumerate}
    \item[(1)] an \textbf{objective-level} re-weighting of the loss, \emph{Item Loss
    Equalization} (Section~\ref{sec:popularity-aware_lightGCN_with_item_loss_equalization}),
    which reshapes per-interaction gradients but leaves the propagation graph itself
    untouched;
    \item[(2)] a \textbf{structure-level} perturbation of the propagation graph, \emph{Degree-Aware
    Graph Augmentation} (Section~\ref{sec:optional_degree_aware_graph_augmentation}), which
    weakens the structural dominance of popular items during message passing but does not
    directly shape the resulting embedding geometry;
    \item[(3)] a \textbf{representation-level} contrastive regularizer computed across two
    augmented views of the graph (Section~\ref{sec:contrastive_learning}), which explicitly
    encourages a well-spread embedding geometry but leaves the negative sampling distribution
    untouched;
    \item[(4)] a \textbf{sampling-level} bias in the negative distribution, \emph{Popularity-Aware
    Negative Sampling} (Section~\ref{sec:popularity_aware_negative_sampling}), which closes
    this last gap by directly and repeatedly challenging popular items during optimization.
\end{enumerate}
All four components are implemented and independently switchable
(Section~\ref{sec:hyperparameters}); however, the configurations reported in
Section~\ref{sec:results} sweep components (1), (3), and (4), while component (2) is exercised
only in its role as the view-construction mechanism for component (3), rather than as a
stand-alone main-pass regularizer --- see the note accompanying
Table~\ref{tab:operating-points} for details.
Section~\ref{sec:joint_training_objective} combines the four components into a single
training objective and summarizes the full training procedure, and
Section~\ref{sec:model_selection_protocol} describes the protocol used to select and evaluate
the final model without leaking test-set information into training decisions.

\subsubsection{Backbone and Item Loss Equalization}
\label{sec:popularity-aware_lightGCN_with_item_loss_equalization}

We first describe the LightGCN backbone and its base training objective, which every
subsequent component in this section builds upon, before introducing our first
bias-mitigation component, Item Loss Equalization.

\paragraph{Backbone.}
Let $\mathcal{U}$ and $\mathcal{I}$ denote the sets of users and items, and let
$e_u^{(0)}, e_i^{(0)} \in \mathbb{R}^d$ be their layer-0 (ego) embeddings, held in a single
embedding table indexed over $\mathcal{U}\cup\mathcal{I}$. Following LightGCN, we propagate
embeddings over the user--item bipartite graph using the symmetric normalized adjacency
$\tilde{A} = D^{-1/2} A D^{-1/2}$, without feature transformation or non-linear activation:
\begin{equation}
    E^{(k+1)} = \tilde{A}\, E^{(k)}, \qquad
    E = \frac{1}{K+1}\sum_{k=0}^{K} E^{(k)},
    \label{eq:lightgcn-propagation}
\end{equation}
where $E^{(k)}$ stacks the embeddings of all nodes at layer $k$ and $K$ is the number of
propagation layers. The final user and item embeddings $e_u, e_i$ are read off from $E$, and
the predicted preference score is their inner product, $s(u,i) = e_u^{\top} e_i$.

\paragraph{Base objective.}
Positive pairs $(u,i^+)$ are sampled uniformly from the set of training interactions, and the
backbone is trained with the Bayesian Personalized Ranking (BPR) loss~\cite{rendle2009bpr}
over the resulting triplets $(u,i^+,i^-)$, where $i^-$ is a sampled negative
(Section~\ref{sec:popularity_aware_negative_sampling} describes our negative sampler):
\begin{equation}
    \ell(u,i^+,i^-) = -\log\sigma\big(s(u,i^+) - s(u,i^-)\big), \qquad
    \mathcal{L}_{BPR} = \frac{1}{|\mathcal{B}|}\sum_{(u,i^+,i^-)\in\mathcal{B}} \ell(u,i^+,i^-),
    \label{eq:bpr}
\end{equation}
together with an $\ell_2$ penalty on the ego embeddings of the nodes present in the batch,
averaged over the batch:
\begin{equation}
    \mathcal{L}_{reg} = \frac{1}{|\mathcal{B}|}\sum_{(u,i^+,i^-)\in\mathcal{B}}
    \Big(\lVert e_u^{(0)}\rVert_2^2 + \lVert e_{i^+}^{(0)}\rVert_2^2 + \lVert e_{i^-}^{(0)}\rVert_2^2\Big),
    \label{eq:l2reg}
\end{equation}
weighted by $\lambda_{reg}$ in the total objective (Eq.~\eqref{eq:total-loss}).

\paragraph{Item Loss Equalization (ILE).}
To counteract popularity bias directly at the objective level, we partition the item catalog
into popularity groups --- \emph{head}, \emph{middle}, \emph{tail} --- based on percentiles
of item degree computed on the training graph (tail: below the 50th percentile; head: above
the 80th percentile). For a training mini-batch $\mathcal{B}$, let $\mathcal{B}_{head}$ and
$\mathcal{B}_{tail}$ be the subsets of triplets whose positive item $i^+$ falls in the head
and tail groups, respectively (triplets whose $i^+$ falls in the middle group belong to
neither subset and do not contribute to $\mathcal{L}_{ILE}$), and let
\begin{equation}
    \bar\ell_{g} = \frac{1}{|\mathcal{B}_{g}|}\sum_{(u,i^+,i^-)\in \mathcal{B}_{g}} \ell(u,i^+,i^-),
    \qquad g \in \{head, tail\},
\end{equation}
be the mean BPR loss within each group (the penalty is only evaluated when both groups are
non-empty in the batch; otherwise it is defined as zero). We define the Item Loss
Equalization penalty as the squared discrepancy between the two group losses,
\begin{equation}
    \mathcal{L}_{ILE} = \big(\bar\ell_{head} - \bar\ell_{tail}\big)^2 .
    \label{eq:ile}
\end{equation}
Because tail items receive fewer positive interactions during training, $\bar\ell_{tail}$ is
typically larger than $\bar\ell_{head}$ under plain BPR training; minimizing
Eq.~\eqref{eq:ile} therefore pulls $\bar\ell_{tail}$ toward $\bar\ell_{head}$, effectively
re-weighting the optimization in favor of tail interactions. The squared, symmetric
formulation is lower-bounded by zero and admits no incentive to instead \emph{increase}
$\bar\ell_{tail}$ to close the gap --- a degenerate failure mode we observed with the naive,
unsigned difference $\bar\ell_{head} - \bar\ell_{tail}$, which collapses the model toward a
popularity-only ranking as its weight increases. ILE is a purely objective-level
intervention: it reshapes the gradient contributed by each interaction according to the
popularity of its positive item but leaves the propagation graph in
Eq.~\eqref{eq:lightgcn-propagation} completely unchanged, isolating its effect from the
structural intervention introduced next.

\subsubsection{Degree-Aware Graph Augmentation}
\label{sec:optional_degree_aware_graph_augmentation}

ILE re-weights the training objective but does not address the structural cause of
popularity bias: high-degree item nodes still dominate neighborhood aggregation in
Eq.~\eqref{eq:lightgcn-propagation}. We introduce an optional, complementary intervention
that perturbs the propagation graph itself, so that popular items contribute a
proportionally smaller share of the aggregated signal during training.

\paragraph{Degree-aware edge dropout.}
For each item $i$ with training-graph degree $d_i$, we define a drop probability that grows
logarithmically with popularity,
\begin{equation}
    p_i = p_{min} + (p_{max} - p_{min}) \cdot \frac{\log(1+d_i)}{\log(1+d_{max})},
    \qquad p_i \in [p_{min}, p_{max}],
    \label{eq:degree-aware-dropout-prob}
\end{equation}
where $d_{max} = \max_i d_i$ and $p_{min}, p_{max}$ are hyperparameters. Popular items are
thus dropped from the propagation graph more aggressively than long-tail items.

\paragraph{Symmetric dropout scheme.}
The bipartite graph is stored as an undirected edge list, i.e., every training interaction
$(u,i)$ contributes both a $u\!\to\!i$ and an $i\!\to\!u$ edge. To preserve the symmetry
required by the normalized adjacency $\tilde{A}$, the drop decision is sampled once per
interaction --- retaining the $u\!\to\!i$ edge independently with probability $1-p_i$ --- and
mirrored onto the reverse $i\!\to\!u$ edge, rather than drawing the two directions
independently. This yields an augmented adjacency $\hat{A}$ that remains a valid undirected
graph. As an ablation control, we also implement \emph{uniform} edge dropout, which drops
every edge independently with a constant probability regardless of item degree, isolating
the contribution of the degree-aware weighting in Eq.~\eqref{eq:degree-aware-dropout-prob}
from that of dropout regularization in general. We refer to the degree-aware/uniform choice
as the \emph{dropout scheme}, reused identically by the contrastive branch in
Section~\ref{sec:contrastive_learning}.

\paragraph{Usage.}
This augmentation is optional in two independent senses: (i) it can be toggled on or off as a
stand-alone regularizer applied to the forward pass used to compute $\mathcal{L}_{BPR}$ and
$\mathcal{L}_{ILE}$, and (ii) it is, independently of (i), also used to construct the two
augmented views consumed by the contrastive objective described next. At inference time the
augmentation is always disabled: item ranking is computed on the original, unperturbed graph,
so the augmentation acts purely as a training-time regularizer. In the configurations
reported in Section~\ref{sec:results}, usage (ii) is always active while usage (i) is not: the
main forward pass is never itself augmented, so the structural intervention studied here is
validated only through the contrastive branch of Section~\ref{sec:contrastive_learning}.

\subsubsection{Contrastive Regularization across Augmented Views}
\label{sec:contrastive_learning}

Degree-aware augmentation (Section~\ref{sec:optional_degree_aware_graph_augmentation})
weakens the structural dominance of popular items during message passing, but it does not
directly control the geometry of the resulting embedding space: item embeddings can still
collapse into a narrow region dominated by popular items even on an augmented graph. We
therefore add a representation-level regularizer that explicitly targets this geometry.
Following the self-supervised graph learning paradigm~\cite{wu2021sgl}, we add an auxiliary
contrastive objective that encourages the embedding space to remain well-spread rather than
collapsing toward a small set of popular items. At each training step we sample two
independent augmented graphs $\hat{A}^{(1)}, \hat{A}^{(2)}$ from the \emph{original} graph
$A$, using the same edge-dropout scheme introduced in
Section~\ref{sec:optional_degree_aware_graph_augmentation} --- degree-aware
(Eq.~\eqref{eq:degree-aware-dropout-prob}) in our main experiments, or uniform in the
ablation control --- independently of whether the main forward pass used for
Eq.~\eqref{eq:bpr} is itself augmented. Each view is propagated through
Eq.~\eqref{eq:lightgcn-propagation} to obtain node embeddings $E^{(1)}, E^{(2)}$, and the two
views of the same node are treated as a positive pair, with all other nodes in the batch
serving as negatives (InfoNCE):
\begin{equation}
    \mathcal{L}_{NCE}(\mathcal{N}) = -\frac{1}{|\mathcal{N}|}\sum_{n\in\mathcal{N}}
    \log \frac{\exp\!\big(\hat{z}_n^{(1)\top}\hat{z}_n^{(2)}/\tau\big)}
    {\sum_{m\in\mathcal{N}} \exp\!\big(\hat{z}_n^{(1)\top}\hat{z}_m^{(2)}/\tau\big)},
    \label{eq:infonce}
\end{equation}
where $\hat{z}_n^{(v)} = E^{(v)}_n / \lVert E^{(v)}_n \rVert$ is the $\ell_2$-normalized
embedding of node $n$ in view $v$, and $\tau$ is a temperature hyperparameter. We apply
Eq.~\eqref{eq:infonce} separately to the set of \emph{unique} users and the set of unique
positive items present in the mini-batch, $\mathcal{U}_{\mathcal{B}}$ and
$\mathcal{I}^{+}_{\mathcal{B}}$, and sum the two terms:
\begin{equation}
    \mathcal{L}_{CL} = \mathcal{L}_{NCE}(\mathcal{U}_{\mathcal{B}})
                      + \mathcal{L}_{NCE}(\mathcal{I}^{+}_{\mathcal{B}}).
\end{equation}
The InfoNCE objective is known to jointly promote \emph{alignment} of positive pairs and
approximately \emph{uniform} coverage of the embedding hypersphere~\cite{wang2020alignuniform};
the latter property directly discourages the embedding collapse toward a small set of popular
items observed under plain BPR training, and is expected to increase catalog coverage.

\subsubsection{Popularity-Aware Negative Sampling}
\label{sec:popularity_aware_negative_sampling}

The three components above reshape the training objective, the propagation graph, and the
embedding geometry, respectively, but none of them alters \emph{which} items are used as
negatives during optimization: negatives are still drawn uniformly, so popular items are not
disproportionately challenged even though they are precisely the items that most need to be
pushed down in the ranking. We close this last gap with a sampling-level intervention.
Standard BPR training draws the negative item $i^-$ uniformly from the catalog. We instead
sample negatives from a degree-biased categorical distribution,
\begin{equation}
    i^- \sim \mathrm{Categorical}(w), \qquad w_i \propto d_i^{\beta}, \qquad \beta \ge 0,
    \label{eq:pop-aware-negative-sampling}
\end{equation}
where $\beta$ controls the sharpness of the bias toward popular items ($\beta=0$ recovers
uniform sampling). Sampled negatives that coincide with an item the user has already
interacted with in training are rejected and re-drawn (subject to a bounded number of
resampling attempts), preserving the standard implicit-BPR assumption that $i^-$ is a
non-interacted item. Because popular items are sampled as negatives far more frequently under
Eq.~\eqref{eq:pop-aware-negative-sampling}, their score is discouraged more often and more
strongly during training, directly counteracting their tendency to dominate the top-$K$
ranking irrespective of true user preference.

\subsubsection{Joint Training Objective}
\label{sec:joint_training_objective}

The full training objective for a mini-batch combines all four components,
\begin{equation}
    \mathcal{L} = \mathcal{L}_{BPR} + \lambda_{reg}\,\mathcal{L}_{reg}
                + \lambda_{ILE}\, \mathcal{L}_{ILE} + \lambda_{CL}\, \mathcal{L}_{CL},
    \label{eq:total-loss}
\end{equation}
where $\mathcal{L}_{BPR}$ (Eq.~\eqref{eq:bpr}) is computed with the popularity-aware negative
sampler of Eq.~\eqref{eq:pop-aware-negative-sampling}, optionally over the augmented graph of
Eq.~\eqref{eq:degree-aware-dropout-prob}; $\mathcal{L}_{ILE}$ reuses the same per-triplet
losses from that forward pass, grouped by the popularity of $i^+$ (Eq.~\eqref{eq:ile}); and
$\mathcal{L}_{CL}$ is computed on two freshly sampled augmented views of the original graph
(Eq.~\eqref{eq:infonce}). Each of $\lambda_{ILE}$, $\lambda_{CL}$, the negative-sampling
sharpness $\beta$, and the two graph-side switches (main-pass augmentation and contrastive
learning) can be set to zero independently. The ablation reported in Section~\ref{sec:results}
sweeps $\lambda_{ILE}$, $\lambda_{CL}$, and $\beta$, while keeping the main-pass augmentation
switch off throughout (Table~\ref{tab:operating-points}); a standalone ablation of main-pass
augmentation is left for future work. At inference time,
item ranking for a user $u$ is computed with the score $s(u,i) = e_u^\top e_i$ evaluated on
the original, unaugmented graph: neither the augmentation nor the contrastive branch is
active at inference, so both act purely as training-time regularizers on an otherwise
unmodified LightGCN scoring function. Algorithm~\ref{alg:training} summarizes one training
epoch.

\begin{algorithm}[t]
\caption{One training epoch of Popularity-Aware LightGCN}
\label{alg:training}
\begin{algorithmic}[1]
\REQUIRE training graph $A$, item degrees $\{d_i\}$, popularity groups, weights
$\lambda_{ILE}, \lambda_{CL}$, temperature $\tau$, sampling sharpness $\beta$, switches for
main-pass augmentation (AUG) and contrastive learning (CL)
\FOR{each mini-batch step}
    \STATE Sample positive pairs $(u,i^+)$ uniformly from the training interactions
    \STATE Sample negatives $i^- \sim \mathrm{Categorical}(d^{\beta})$
    (Eq.~\eqref{eq:pop-aware-negative-sampling}); reject and re-draw if $i^-$ is a training
    positive of $u$
    \STATE $\hat{A}_{\mathrm{main}} \leftarrow$ degree-aware (or uniform) dropout of $A$ if
    AUG is enabled, else $A$
    \STATE Propagate $\hat{A}_{\mathrm{main}}$ (Eq.~\eqref{eq:lightgcn-propagation}) to obtain
    $s(u,i^+)$, $s(u,i^-)$, and the ego embeddings of the batch
    \STATE Compute $\mathcal{L}_{BPR}$ (Eq.~\eqref{eq:bpr}) and $\mathcal{L}_{reg}$ (Eq.~\eqref{eq:l2reg})
    \IF{ILE is enabled}
        \STATE Partition the batch by the popularity group of $i^+$; compute
        $\mathcal{L}_{ILE}$ (Eq.~\eqref{eq:ile})
    \ENDIF
    \IF{CL is enabled}
        \STATE Sample two fresh views $\hat{A}^{(1)}, \hat{A}^{(2)}$ of the original graph $A$
        via the configured edge-dropout scheme; propagate both
        \STATE Compute $\mathcal{L}_{CL}$ on the batch's unique users and positive items
        (Eq.~\eqref{eq:infonce})
    \ENDIF
    \STATE Assemble $\mathcal{L}$ (Eq.~\eqref{eq:total-loss}); backpropagate and update
    parameters with Adam
\ENDFOR
\STATE Every $E$ epochs, evaluate Recall@$K$ on the validation split and retain the checkpoint
with the highest validation Recall@$K$; stop early after $P$ evaluations without improvement
\end{algorithmic}
\end{algorithm}

\subsubsection{Model Selection and Evaluation Protocol}
\label{sec:model_selection_protocol}

Having described the model and its unified training objective, we conclude with the protocol
used to select and evaluate the trained model, which underpins the validity of the results
reported in Section~\ref{sec:results}. To obtain an unbiased estimate of generalization
performance, the validation and test splits
are strictly separated in role throughout training. Every $E$ epochs, the current model is
evaluated on the validation split and the checkpoint achieving the highest validation
Recall@$K$ is retained; training stops early if no improvement is observed for $P$ consecutive
evaluations (Algorithm~\ref{alg:training}). The test split is touched \emph{exactly once},
after training has concluded, using the retained best-validation checkpoint; at no point is a
test-set metric used to make a training or model-selection decision. All test-set results
reported in Section~\ref{sec:results} therefore reflect genuine out-of-sample generalization
rather than a checkpoint selected in hindsight on the test set. Unless otherwise noted,
$K=20$ throughout: Recall@20 is used both as the model-selection criterion in
Algorithm~\ref{alg:training} and as the primary reported accuracy metric in Section~\ref{sec:results}.

\subsection{Implementation Details}
\label{sec:implementation_details}

\paragraph{Architecture and initialization.}
All models share an embedding dimension of $d=64$. The backbone
(Section~\ref{sec:popularity-aware_lightGCN_with_item_loss_equalization}) uses $K=2$
propagation layers in our main experiments, which we found to avoid the over-smoothing effect
observed with a deeper, 3-layer backbone on this dataset (Section~\ref{sec:results}). Embedding
weights are initialized from $\mathcal{N}(0, 0.1^2)$.

\paragraph{Popularity groups.}
Item popularity groups (head/middle/tail,
Section~\ref{sec:popularity-aware_lightGCN_with_item_loss_equalization}) are computed once
from training-graph item degrees using the 50th and 80th percentile cutoffs described above,
and held fixed throughout training and evaluation.

\paragraph{Optimization.}
All models are trained with Adam using a learning rate of $10^{-3}$, an $\ell_2$ weight of
$\lambda_{reg}=10^{-4}$ (Eq.~\eqref{eq:l2reg}), and mini-batches of 4096 triplets. Gradients
are clipped to a maximum global norm of $1.0$ after every backward pass. Models are trained
for up to 100 epochs under the early-stopping and checkpointing protocol of
Section~\ref{sec:model_selection_protocol} ($E=5$, $P=10$).

The complete list of hyperparameters and their explored ranges is given in
Section~\ref{sec:hyperparameters} (Table~\ref{tab:hyperparams-master}), together with the
exact combination selected for each configuration reported in Section~\ref{sec:results}
(Table~\ref{tab:operating-points}).

\paragraph{Software, hardware, and reproducibility.}
The model is implemented in PyTorch and trained on a single NVIDIA A100 GPU. Each reported
configuration is trained with three random seeds (42, 0, 1); the model parameters and the
mini-batch sampling stream are seeded separately, so that a given run is exactly
reproducible. Each run periodically checkpoints its model and optimizer state, allowing
training to be resumed after interruption without altering the final result, and logs
per-epoch training and validation metrics for post-hoc inspection.

\subsection{Hyperparameters}
\label{sec:hyperparameters}

Table~\ref{tab:hyperparams-master} lists every hyperparameter introduced in
Section~\ref{sec:proposed_method} and Section~\ref{sec:implementation_details}, together with
the value or range used in this work. Table~\ref{tab:operating-points} then gives the exact
combination selected for each named configuration reported in Section~\ref{sec:results}.

\begin{table}[t]
\centering
\caption{All hyperparameters of Popularity-Aware LightGCN.}
\label{tab:hyperparams-master}
\fitwidthtable{%
\begin{tabular}{lll}
\toprule
Symbol & Description & Value / range \\
\midrule
$d$ & Embedding dimension & $64$ \\
$K$ & Propagation layers (Eq.~\eqref{eq:lightgcn-propagation}) & $\{2, 3\}$ \\
--- & Adam learning rate & $10^{-3}$ \\
$\lambda_{reg}$ & $\ell_2$ weight (Eq.~\eqref{eq:l2reg}) & $10^{-4}$ \\
--- & Mini-batch size & $4096$ triplets \\
--- & Gradient-clipping max norm & $1.0$ \\
--- & Maximum training epochs & $100$ \\
$E$ & Validation frequency (Algorithm~\ref{alg:training}) & every $5$ epochs \\
$P$ & Early-stopping patience (Algorithm~\ref{alg:training}) & $10$ evaluations \\
--- & Head / tail popularity percentile cutoffs & $80$th / $50$th \\
$\lambda_{ILE}$ & ILE weight (Eq.~\eqref{eq:total-loss}) & $\{0, 0.1, 0.5, 1.0\}$ \\
$\lambda_{CL}$ & Contrastive weight (Eq.~\eqref{eq:total-loss}) & $\{0, 0.1, 0.5\}$ \\
$\tau$ & InfoNCE temperature (Eq.~\eqref{eq:infonce}) & $0.2$ \\
$p_{min}, p_{max}$ & Degree-aware dropout range (Eq.~\eqref{eq:degree-aware-dropout-prob}) & $0.1$, $0.4$ \\
$p_{uniform}$ & Uniform-dropout ablation probability & $0.1$ \\
$\beta$ & Negative-sampling sharpness (Eq.~\eqref{eq:pop-aware-negative-sampling}) & $\{0, 0.5\}$ \\
\bottomrule
\end{tabular}}
\end{table}

\begin{table}[t]
\centering
\caption{Exact hyperparameter combination for each configuration reported in
Section~\ref{sec:results} (test-set results, mean over 3 seeds: 42, 0, 1). All configurations use
$d=64$, $K=2$, $\tau=0.2$, and the degree-aware dropout scheme; the main forward pass used for
$\mathcal{L}_{BPR}$ and $\mathcal{L}_{ILE}$ is \emph{not} itself augmented in these results
(aug\_main~$=$~off for every row). Consequently, Degree-Aware Graph Augmentation
(Section~\ref{sec:optional_degree_aware_graph_augmentation}) is validated here only
indirectly, through the view construction of the contrastive branch
(Section~\ref{sec:contrastive_learning}); a standalone ablation of main-pass augmentation is
left for future work.}
\label{tab:operating-points}
\fitwidthtable{%
\begin{tabular}{lccc}
\toprule
Configuration & $\lambda_{ILE}$ & $\lambda_{CL}$ & $\beta$ \\
\midrule
LightGCN (baseline) & $0$ & $0$ & $0$ \\
PopAware --- accuracy & $0.1$ & $0.1$ & $0.5$ \\
PopAware-BEST --- balanced & $1.0$ & $0.1$ & $0.5$ \\
PopAware --- high-tail & $1.0$ & $0.1$ & $0$ \\
PopAware --- fairness & $1.0$ & $0.5$ & $0$ \\
\bottomrule
\end{tabular}}
\end{table}

\subsection{Results for the Proposed Method}
\label{sec:results}

\paragraph{Setup recap.}
We evaluate on MovieLens-1M under the leave-one-out protocol of
Section~\ref{sec:model_selection_protocol}: $6{,}034$ users, $3{,}533$ items, $563{,}204$
training interactions, and one held-out validation and one held-out test interaction per
user. All results below follow the exact hyperparameter combinations of
Table~\ref{tab:operating-points} and are averaged over three seeds ($42$, $0$, $1$), reported
as mean~$\pm$~standard deviation. We report the accuracy metrics (Recall@10, NDCG@10,
Recall@20, NDCG@20) and the popularity-bias metrics (TailRecall@20, Coverage@20, ARP@20,
Tail/Middle/Head Exposure@20) defined in Section~[Metrics]. Recalling the motivation of
Section~\ref{sec:proposed_method}, a successful outcome requires TailRecall@20,
Coverage@20, TailExposure@20, and MiddleExposure@20 to \emph{increase}, ARP@20 and
HeadExposure@20 to \emph{decrease}, and Recall@20/NDCG@20 to not decrease sharply relative to
the LightGCN backbone.

\paragraph{Main results.}
Table~\ref{tab:results-accuracy} reports accuracy and Table~\ref{tab:results-bias} reports
popularity-bias metrics, for the LightGCN backbone, the four operating points of
Popularity-Aware LightGCN (Table~\ref{tab:operating-points}), and two external reference
methods reported elsewhere in this work, MostPopular and BPR-MF (single run each, no seed
variance available). \emph{PopAware-BEST} improves \emph{every} metric relative to the
LightGCN backbone: Recall@20 and NDCG@20 both increase (rather than merely holding steady),
while TailRecall@20 rises by a factor of $6.5\times$, Coverage@20 by $35\%$, and ARP@20 and
HeadExposure@20 fall by $6.3\%$ and $7.9\%$, respectively. The remaining three operating
points trace out a frontier along the accuracy--fairness trade-off: \emph{accuracy} attains
the highest Recall@20/NDCG@20 in the table, while \emph{high-tail} attains the highest
Recall@10/NDCG@10 and the highest TailRecall@20; and \emph{fairness} attains the best
Coverage@20, ARP@20, TailExposure@20, MiddleExposure@20, and HeadExposure@20 of all seven
methods, at the cost of the largest accuracy drop. No baseline or external reference method
attains the best value on any metric in either table --- every bolded entry belongs to one of
the four PopAware operating points.

\begin{table}[t]
\centering
\footnotesize
\caption{Accuracy metrics on the MovieLens-1M test set. PopAware rows report mean $\pm$
std over 3 seeds; MostPopular and BPR-MF are single-run external references.}
\label{tab:results-accuracy}
\fitwidthtable{%
\begin{tabular}{lcccc}
\toprule
Method & Recall@10$\uparrow$ & NDCG@10$\uparrow$ & Recall@20$\uparrow$ & NDCG@20$\uparrow$ \\
\midrule
MostPopular (ref.) & $0.0403$ & $0.0189$ & $0.0723$ & $0.0269$ \\
BPR-MF (ref.) & $0.0742$ & $0.0369$ & $0.1309$ & $0.0511$ \\
LightGCN (backbone) & $0.0725\pm0.0013$ & $0.0356\pm0.0002$ & $0.1287\pm0.0005$ & $0.0497\pm0.0002$ \\
PopAware --- accuracy & $0.0771\pm0.0009$ & $0.0378\pm0.0003$ & $\mathbf{0.1367\pm0.0035}$ & $\mathbf{0.0527\pm0.0006}$ \\
PopAware-BEST --- balanced & $0.0771\pm0.0012$ & $0.0376\pm0.0003$ & $0.1338\pm0.0030$ & $0.0518\pm0.0004$ \\
PopAware --- high-tail & $\mathbf{0.0783\pm0.0006}$ & $\mathbf{0.0382\pm0.0002}$ & $0.1352\pm0.0020$ & $0.0524\pm0.0004$ \\
PopAware --- fairness & $0.0638\pm0.0020$ & $0.0310\pm0.0005$ & $0.1052\pm0.0039$ & $0.0415\pm0.0011$ \\
\bottomrule
\end{tabular}}
\end{table}

\begin{table}[t]
\centering
\footnotesize
\caption{Popularity-bias metrics on the MovieLens-1M test set. Arrows indicate the direction
of improvement. PopAware rows report mean $\pm$ std over 3 seeds; MostPopular and BPR-MF are
single-run external references.}
\label{tab:results-bias}
\fitwidthtable{%
\begin{tabular}{lcccccc}
\toprule
Method & TailRec.@20$\uparrow$ & Cov.@20$\uparrow$ & ARP@20$\downarrow$ & TailExp.@20$\uparrow$ & MidExp.@20$\uparrow$ & HeadExp.@20$\downarrow$ \\
\midrule
MostPopular (ref.) & $0.0000$ & $0.0498$ & $7.530$ & $0.0000$ & $0.0000$ & $1.0000$ \\
BPR-MF (ref.) & $0.0109$ & $0.628$ & $6.518$ & $0.0086$ & $0.1140$ & $0.8774$ \\
LightGCN (backbone) & $0.0050\pm0.0009$ & $0.398\pm0.004$ & $6.833\pm0.009$ & $0.0015\pm0.0001$ & $0.0398\pm0.0006$ & $0.9587\pm0.0006$ \\
PopAware --- accuracy & $0.0143\pm0.0023$ & $0.462\pm0.029$ & $6.500\pm0.051$ & $0.0097\pm0.0011$ & $0.0732\pm0.0133$ & $0.9171\pm0.0144$ \\
PopAware-BEST --- balanced & $0.0324\pm0.0049$ & $0.538\pm0.029$ & $6.400\pm0.042$ & $0.0320\pm0.0012$ & $0.0854\pm0.0104$ & $0.8826\pm0.0114$ \\
PopAware --- high-tail & $\mathbf{0.0399\pm0.0023}$ & $0.499\pm0.010$ & $6.564\pm0.008$ & $0.0363\pm0.0014$ & $0.0521\pm0.0019$ & $0.9117\pm0.0005$ \\
PopAware --- fairness & $0.0274\pm0.0035$ & $\mathbf{0.713\pm0.003}$ & $\mathbf{6.138\pm0.001}$ & $\mathbf{0.0482\pm0.0006}$ & $\mathbf{0.1184\pm0.0013}$ & $\mathbf{0.8333\pm0.0016}$ \\
\bottomrule
\end{tabular}}
\end{table}

\paragraph{Effect of the individual hyperparameters.}
The four configurations of Table~\ref{tab:operating-points} vary $\lambda_{ILE}$,
$\lambda_{CL}$, and $\beta$ \emph{jointly}, so no pair of them isolates a single factor. To
isolate each factor individually, we additionally consult the single-seed grid search used to
select these operating points (all at $K=2$), changing exactly one hyperparameter at a time
from a common base setting. Table~\ref{tab:results-sweep-isolated} reports the resulting
Recall@20, TailRecall@20, and Coverage@20.

Increasing $\lambda_{ILE}$ from $0.1$ to $1.0$ (at fixed $\lambda_{CL}=0.1$, $\beta=0$) raises
TailRecall@20 monotonically ($0.0131 \to 0.0262 \to 0.0374$) while leaving Recall@20
essentially unchanged ($0.1356 \to 0.1352 \to 0.1354$), without the collapse toward a
popularity-only ranking observed with the earlier, incorrectly signed formulation of ILE
(Section~\ref{sec:popularity-aware_lightGCN_with_item_loss_equalization}). Increasing
$\lambda_{CL}$ from $0.1$ to $0.5$ (at fixed $\lambda_{ILE}=0.1$, $\beta=0$) trades accuracy
for fairness: Coverage@20 rises by $47\%$ ($0.444\to0.653$), while Recall@20 falls by $18\%$
($0.1356\to0.1112$); this $17$--$19\%$ drop is consistent across all six matched
$(\lambda_{ILE},\beta)$ pairs in the full grid, not just the one shown. Increasing $\beta$
from $0$ to $0.5$ (at fixed $\lambda_{ILE}=1.0$, $\lambda_{CL}=0.1$) raises Coverage@20 by
$11\%$ and lowers HeadExposure@20, at the cost of a $20\%$ \emph{relative} drop in
TailRecall@20 ($0.0374\to0.0299$) alongside a negligible change in Recall@20 --- i.e.,
$\beta$ and $\lambda_{ILE}$ partially compete for the same TailRecall@20 gain rather than
combining additively.

\begin{table}[t]
\centering
\footnotesize
\caption{Isolated effect of each hyperparameter: single-seed grid search at $K=2$, changing
exactly one factor at a time from the base row. Supplementary evidence for hyperparameter
selection, not a substitute for the 3-seed results of Tables~\ref{tab:results-accuracy}--\ref{tab:results-bias}.}
\label{tab:results-sweep-isolated}
\fitwidthtable{%
\begin{tabular}{lccc}
\toprule
Configuration ($K=2$, single seed) & Recall@20$\uparrow$ & TailRecall@20$\uparrow$ & Coverage@20$\uparrow$ \\
\midrule
Base: $\lambda_{ILE}{=}0.1,\ \lambda_{CL}{=}0.1,\ \beta{=}0$ & $0.1356$ & $0.0131$ & $0.444$ \\
$\ \to \lambda_{ILE}{=}0.5$ (others as base) & $0.1352$ & $0.0262$ & $0.498$ \\
$\ \to \lambda_{ILE}{=}1.0$ (others as base) & $0.1354$ & $0.0374$ & $0.503$ \\
$\ \to \lambda_{CL}{=}0.5$ (others as base) & $0.1112$ & $0.0131$ & $0.653$ \\
\midrule
Base$'$: $\lambda_{ILE}{=}1.0,\ \lambda_{CL}{=}0.1,\ \beta{=}0$ & $0.1354$ & $0.0374$ & $0.503$ \\
$\ \to \beta{=}0.5$ (others as Base$'$) & $0.1347$ & $0.0299$ & $0.558$ \\
\bottomrule
\end{tabular}}
\end{table}

Table~\ref{tab:results-layers} compares $K=2$ against $K=3$. For the \emph{plain} LightGCN
backbone (no ILE, CL, or negative-sampling), $K=2$ clearly dominates $K=3$ on every metric,
consistent with over-smoothing of the deeper backbone amplifying popularity bias. This is the
comparison that motivated our choice of $K=2$ for the headline configurations. Once the full
method is active, however, this advantage becomes small and inconsistent in sign: averaged
over the $12$ matched $(\lambda_{ILE},\lambda_{CL},\beta)$ grid points (single seed, both
values of $\beta$), $K=3$ is marginally \emph{higher} than $K=2$ on both Recall@20 (by
$+0.0009$ on average; range $-0.0020$ to $+0.0037$ across the 12 pairs) and Coverage@20 (by
$+0.010$ on average; range $-0.037$ to $+0.047$). In other words, the bias-mitigation
components appear to substantially offset the over-smoothing penalty of the deeper backbone,
though not consistently enough at a single seed to justify switching the headline
configuration to $K=3$; we retain $K=2$ for the 3-seed results above and leave a 3-seed
confirmation of the full method at $K=3$ to future work.

\begin{table}[t]
\centering
\footnotesize
\caption{Effect of the number of propagation layers $K$ on the \emph{plain} LightGCN
backbone (single seed; no ILE, CL, or negative-sampling).}
\label{tab:results-layers}
\fitwidthtable{%
\begin{tabular}{lccccc}
\toprule
Setting & Recall@20$\uparrow$ & TailRecall@20$\uparrow$ & Coverage@20$\uparrow$ & ARP@20$\downarrow$ & HeadExposure@20$\downarrow$ \\
\midrule
LightGCN backbone, $K=2$ & $0.1279$ & $0.0056$ & $0.392$ & $6.846$ & $0.960$ \\
LightGCN backbone, $K=3$ & $0.1206$ & $0.0037$ & $0.344$ & $6.944$ & $0.972$ \\
\bottomrule
\end{tabular}}
\end{table}

\paragraph{Statistical robustness.}
Because PopAware-BEST is trained with the same three seeds as the LightGCN backbone,
Table~\ref{tab:results-significance} checks whether the improvements above exceed the
run-to-run variance rather than reflecting a single fortunate run. For every one of the five
target metrics, the mean~$\pm$~std interval of PopAware-BEST does not overlap with that of the
LightGCN backbone, indicating that the improvement is unlikely to be an artifact of seed
variance.

\begin{table}[t]
\centering
\footnotesize
\caption{Non-overlap check between LightGCN and PopAware-BEST (mean $\pm$ 1 std per side).}
\label{tab:results-significance}
\fitwidthtable{%
\begin{tabular}{lccc}
\toprule
Metric & LightGCN interval & PopAware-BEST interval & Overlap? \\
\midrule
Recall@20 $\uparrow$ & $[0.1282, 0.1292]$ & $[0.1308, 0.1368]$ & No \\
TailRecall@20 $\uparrow$ & $[0.0041, 0.0059]$ & $[0.0275, 0.0373]$ & No \\
Coverage@20 $\uparrow$ & $[0.394, 0.402]$ & $[0.509, 0.567]$ & No \\
ARP@20 $\downarrow$ & $[6.824, 6.842]$ & $[6.358, 6.442]$ & No \\
HeadExposure@20 $\downarrow$ & $[0.9581, 0.9593]$ & $[0.8712, 0.8940]$ & No \\
\bottomrule
\end{tabular}}
\end{table}

\paragraph{Checklist against the stated expectations.}
Having established that the gains of PopAware-BEST over the LightGCN backbone are unlikely to
be seed noise, we now check them systematically against the design goals set out at the start
of Section~\ref{sec:proposed_method}. Table~\ref{tab:results-checklist} maps each criterion
onto the observed change from LightGCN to PopAware-BEST. All six criteria are met; notably,
the accuracy criterion is not merely satisfied but exceeded, since Recall@20 and NDCG@20 both
increase rather than staying flat or decreasing.

\begin{table}[t]
\centering
\footnotesize
\caption{Expectations versus observed outcome (LightGCN $\to$ PopAware-BEST).}
\label{tab:results-checklist}
\fitwidthtable{%
\begin{tabular}{lcc}
\toprule
Expectation & Observed change & Met? \\
\midrule
TailRecall@20 $\uparrow$ & $0.0050 \to 0.0324$ ($6.5\times$) & Yes \\
Coverage@20 $\uparrow$ & $0.398 \to 0.538$ ($+35\%$) & Yes \\
ARP@20 $\downarrow$ & $6.833 \to 6.400$ ($-6.3\%$) & Yes \\
HeadExposure@20 $\downarrow$ & $0.9587 \to 0.8826$ ($-7.9\%$) & Yes \\
Recall@20$\uparrow$ / NDCG@20$\uparrow$, no sharp drop & $+4.0\%$ / $+4.2\%$ & Yes (exceeded) \\
TailRecall@20 ``acceptable'' (neither negligible nor extreme) & $6.5\times$ baseline, below the $8.0\times$ of \emph{high-tail} & Yes \\
\bottomrule
\end{tabular}}
\end{table}

\paragraph{Comparison against the strongest external baseline.}
The checklist above confirms PopAware-BEST improves on its own backbone; we next ask how it
fares against a competitive method it was never designed against. BPR-MF is the strongest of
the two external reference methods on both accuracy and fairness, and is therefore the more
informative comparison point. PopAware-BEST exceeds BPR-MF on
Recall@20 ($0.1338$ vs.\ $0.1309$), NDCG@20 ($0.0518$ vs.\ $0.0511$), TailRecall@20 ($0.0324$
vs.\ $0.0109$, roughly $3\times$), ARP@20 ($6.400$ vs.\ $6.518$), and TailExposure@20
($0.0320$ vs.\ $0.0086$); it trails on Coverage@20 ($0.538$ vs.\ $0.628$) and
MiddleExposure@20 ($0.0854$ vs.\ $0.1140$), and is approximately tied on HeadExposure@20
($0.8826$ vs.\ $0.8774$). The \emph{fairness} operating point instead dominates BPR-MF on
every popularity-bias metric (Coverage@20 $0.713$ vs.\ $0.628$; ARP@20 $6.138$ vs.\ $6.518$;
TailExposure@20 $0.0482$ vs.\ $0.0086$; MiddleExposure@20 $0.1184$ vs.\ $0.1140$;
HeadExposure@20 $0.8333$ vs.\ $0.8774$), at the cost of a lower Recall@20 ($0.1052$ vs.\
$0.1309$). Because BPR-MF is reported from a single run, these comparisons cannot be checked
for statistical significance in the same way as Table~\ref{tab:results-significance}, and
should be read as indicative rather than conclusive.

\paragraph{Supplementary: best-observed-seed results.}
Table~\ref{tab:results-bestseed} additionally reports, for each method, the single seed that
attains the highest Recall@20 --- \emph{not} the mean~$\pm$~std of Table~\ref{tab:results-accuracy}/\ref{tab:results-bias}. Because this is the maximum of three noisy
measurements, it is a systematically optimistic (upward-biased) estimator: repeating the
experiment with more seeds would keep pushing this number up even if the underlying method
were unchanged. Accordingly, seed~$1$ attains the best Recall@20 for four of the five PopAware
configurations, while seed~$42$ attains it for \emph{fairness} --- there is no single ``best''
seed across configurations, which is itself evidence that these figures should not be read as
representative performance. We report Table~\ref{tab:results-bestseed} only for
completeness; all claims and comparisons in this work are based on
Tables~\ref{tab:results-accuracy}--\ref{tab:results-checklist}.

\begin{table}[t]
\centering
\footnotesize
\caption{Best-observed-seed results (illustrative only; see caveat above). Not used for any
claim in this work.}
\label{tab:results-bestseed}
\fitwidthtable{%
\begin{tabular}{lccccc}
\toprule
Method & Seed & Recall@20$\uparrow$ & NDCG@20$\uparrow$ & TailRecall@20$\uparrow$ & Coverage@20$\uparrow$ \\
\midrule
LightGCN (backbone) & 1 & $0.1293$ & $0.0500$ & $0.0037$ & $0.399$ \\
PopAware --- accuracy & 1 & $0.1407$ & $0.0534$ & $0.0112$ & $0.439$ \\
PopAware-BEST --- balanced & 1 & $0.1369$ & $0.0523$ & $0.0393$ & $0.560$ \\
PopAware --- high-tail & 1 & $0.1376$ & $0.0527$ & $0.0393$ & $0.509$ \\
PopAware --- fairness & 42 & $0.1100$ & $0.0430$ & $0.0299$ & $0.714$ \\
\bottomrule
\end{tabular}}
\end{table}

In summary, PopAware-BEST meets every criterion set out in Section~\ref{sec:proposed_method}
with statistically robust margins, and the broader family of four operating points traces a
frontier that is competitive with, and in several respects exceeds, the strongest external
baseline evaluated in this work. The implications, remaining limitations, and directions for
future work are discussed in Section~[Discussion].
